% TOP_PROBLEM: {"id": 3100068, "title": "Is there a prime between n^(2)and (n+1)^(2 )for every n > 0", "category_id": 1, "category": {"id": 1, "name": "number_theory", "display_name": "Number Theory", "description": "Properties of integers, prime numbers, Diophantine equations.", "slug": "number-theory", "order_index": 1, "created_at": "2026-07-31T15:26:25.670Z"}, "status": "partially_solved"}
% FIRST_POSTED: null
% ENTRY_KIND: statement_only
% Imported from: batch11.tex; UnsolvedMath ID 3100068
\documentclass[11pt]{article}
\usepackage{amsmath,amssymb,mathtools}
\usepackage{fontspec,unicode-math}
\setmainfont{Latin Modern Roman}
\setmathfont{Latin Modern Math}
\usepackage{xurl,hyperref}
\newcommand{\sref}[2]{\hyperlink{source:#1:#2}{[S#2]}}
\newcommand{\cataloguescope}[1]{\paragraph{Mathematical statement.}#1}
\begin{document}
% UnsolvedMath ID 3100068; source catalogue code AMR-030-0068
\setcounter{section}{3100067}
\section{Is there a prime between n\textasciicircum{}(2)and (n+1)\textasciicircum{}(2 )for every n > 0}\label{3100068}
{\small\sffamily UnsolvedMath ID 3100068 · Number Theory}
\subsection{Definitions and mathematical statement}

\paragraph{Shared notation.}$\mathbb N=\{1,2,\ldots\}$, and $\mathbb Z,\mathbb Q,\mathbb R,\mathbb C$ denote the integers, rationals, reals and complex numbers. Any other conventions are specified locally.


A prime is an integer $p\ge2$ whose only positive divisors are $1$ and $p$. The interval between consecutive squares is taken to be open.

\paragraph{Statement recorded in the supplied source.}
For every $n\in\mathbb N$, does there exist a prime $p$ such that $n^2<p<(n+1)^2$? \sref{3100068}{2}

\subsection{Short English statement}
Is there always a prime strictly between the squares of two consecutive positive integers?
\subsection{Sources}
\begin{description}
\item[\hypertarget{source:3100068:1}{S1}] ULAM.AI and the UnsolvedMath Contributors, \emph{UnsolvedMath: A Curated Collection of Open Mathematics Problems}, record 3100068 (AMR-030-0068). CC BY 4.0. \url{https://huggingface.co/datasets/ulamai/UnsolvedMath}.
\item[\hypertarget{source:3100068:2}{S2}] Jonathan Sorenson and Jonathan Webster, \emph{An algorithm and computation to verify Legendre\textquotesingle s conjecture up to $3.33\cdot10^{13}$} (2024), Abstract and Section1, Legendre’s conjecture, page1. \url{https://arxiv.org/abs/2401.13753}.
\end{description}
\label{end:3100068}
\end{document}
